% Our Approach -- needs \usepackage{amsmath,amssymb,algorithm,algorithmic}
% \newcommand{\ours}{\textsc{Push+Margin}}   % put in the preamble (rename freely)

\section{Our Approach}
\label{sec:approach}

SemStamp treats every sentence position the same way: it resamples until a sentence lands in the mask and clears the
margin. Whether this can succeed depends on how many different meanings the model is willing to express at that
position, i.e.\ on its \emph{semantic uncertainty}. Where the model is certain, its samples are near-copies of one
another, fall into the same region, and every further sample repeats the same failure until the budget runs out. Where
the model is uncertain, a few samples already cover several regions. Our approach, \ours, makes the watermarking effort
follow the model's uncertainty at two levels: at the sentence level, a \emph{bounded selection} that stops as soon as a
usable sentence is found and gives up cheaply where none exists; at the token level, a \emph{push} that steers word
choices toward the mask only when the next-token distribution has high entropy. The encoder, the key, the mask rule and
the detector are SemStamp's, unchanged.

\subsection{Setup}
\label{sec:setup}

A language model $p_\theta$ writes a text $y = s_1 s_2 \dots s_n$ of sentences in response to a prompt $x$. A sentence
encoder maps a sentence $s$ to a unit vector $e(s) \in \mathbb{R}^D$. The secret key $k$ defines $d$ random hyperplanes
with unit normals $w_1, \dots, w_d$, which split the embedding space into $2^d$ regions; the region $r(s)$ of a sentence
is given by the signs of $\langle e(s), w_j \rangle$. The key also maps the region of the previous sentence to a mask
$M_t = \mathcal{M}_k\big(r(s_{t-1})\big)$, a set of $\gamma 2^d$ regions. For a region $b$ let $\sigma_b \in \{\pm 1\}^d$
denote the side of each hyperplane on which $b$ lies. We measure how far a sentence lies inside the mask by its
\emph{depth}
\begin{equation}
  \delta(s; M) \;=\; \max_{b \in M} \; \min_{1 \le j \le d} \; \sigma_{b,j} \, \langle e(s), w_j \rangle .
  \label{eq:depth}
\end{equation}
$\delta(s; M) > 0$ if and only if $s$ lies in a region of the mask, and then $\delta$ is the cosine distance to the
nearest face of that region. A sentence is \emph{usable} if $\delta(s; M) \ge m$, i.e.\ it is in the mask and clears
SemStamp's margin $m$. We use SemStamp's values $d = 3$, $\gamma = 0.25$ and $m = 0.02$.

\subsection{Watermark capacity and uncertainty}
\label{sec:capacity}

We distinguish two kinds of uncertainty. \emph{Lexical} uncertainty is the entropy of the next-token distribution,
\begin{equation}
  H_t \;=\; -\sum_{v} p_\theta(v \mid x, y_{<t}) \log p_\theta(v \mid x, y_{<t}) ,
  \label{eq:token-entropy}
\end{equation}
which is high when many tokens are plausible continuations. \emph{Semantic} uncertainty is uncertainty about what the
next sentence means \citep{kuhn2023semantic}. At a sentence position we estimate it by drawing $K$ candidate sentences and
computing the entropy of their distribution over the $2^d$ regions,
\begin{equation}
  H_{\text{bin}} \;=\; -\sum_{b=1}^{2^d} \hat q_b \log \hat q_b ,
  \qquad \hat q_b = \tfrac{1}{K}\,\big|\{ i : r(c_i) = b \}\big| ,
  \label{eq:bin-entropy}
\end{equation}
which is $0$ when all candidates share one region and at most $d \ln 2$ when they are spread evenly. $H_{\text{bin}}$ is
a coarse, key-dependent analogue of semantic entropy in which the regions play the role of meaning clusters.

The \emph{capacity} of a position is the probability that at least one of its candidates is usable. If draws were
independent with per-draw success probability $\rho_t$, $K$ draws would succeed with probability $1 - (1 - \rho_t)^K$,
and enough resampling would always succeed. Draws are not independent in meaning: when $H_{\text{bin}} = 0$ every
candidate lands in the same region, so the position can carry the mark only if that one region happens to be in the
mask, and resampling cannot help. Capacity is therefore governed by semantic uncertainty, not by the sampling budget.
This motivates three design principles:
\begin{enumerate}
  \item[(P1)] \textbf{Bound the effort per sentence.} Positions without semantic freedom should be passed over unmarked
  at a small, fixed cost instead of exhausting a large rejection budget.
  \item[(P2)] \textbf{Raise the per-draw success only where the model is uncertain.} Steering a token toward the mask is
  cheap when several tokens are plausible ($H_t$ high) and harmful when one token is (near-)certain, e.g.\ inside a
  name, a number or a fixed phrase.
  \item[(P3)] \textbf{Never change where a sentence or the text ends.} The sentence boundaries define the units that
  the detector scores, and the length of the answer is part of its quality.
\end{enumerate}

\subsection{Bounded selection with a margin stopping rule}
\label{sec:selection}

The first sentence $s_1$ is a single ordinary sample: it carries no mark and only seeds the chain of masks. For every
later position $t$, \ours{} draws candidate sentences in blocks of $B$, at most $K$ in total, each written with the push
of Section~\ref{sec:push}. After each block the candidates are embedded. The search stops at the first block that
contains a usable candidate, and the usable candidate of largest depth is kept. If $K$ candidates have been drawn and none
is usable, the candidate of largest depth is kept anyway; it may lie in the mask without clearing the margin, or outside
the mask, in which case the sentence carries no mark and generation continues (P1). Algorithm~\ref{alg:ours} gives the
full procedure; we use $B = 4$ and $K = 16$.

SemStamp requires every kept sentence to clear the margin and rejects it otherwise. \ours{} uses the margin only as a
\emph{stopping rule}: a block ends the search only if it holds a candidate that is both in the mask and clear of the
margin, so the margin decides how long to search, not what may be kept. Keeping the deepest usable candidate adds
distance from the region's faces at no extra cost, which helps the sentence stay in its region after the text is edited
and re-embedded. Every kept sentence is one of the model's own samples.

\begin{algorithm}[t]
\caption{\ours{}: generation of one watermarked text}
\label{alg:ours}
\begin{algorithmic}[1]
\REQUIRE prompt $x$, key $k$, block size $B$, cap $K$, margin $m$
\STATE $s_1 \sim p_\theta(\cdot \mid x)$ \COMMENT{plain sample, seeds the chain}
\STATE $t \gets 1$
\WHILE{the model does not end the answer after $s_{1:t}$}
  \STATE $M \gets \mathcal{M}_k\big(r(s_t)\big)$; \ $\mathcal{C} \gets \emptyset$
  \REPEAT
    \STATE draw $B$ sentences $c \sim \tilde p_\theta(\cdot \mid x, s_{1:t}; M)$ and add them to $\mathcal{C}$ \COMMENT{push, Eq.~\eqref{eq:tilt}}
  \UNTIL{a candidate of the last block has $\delta(c; M) \ge m$ \OR $|\mathcal{C}| = K$}
  \STATE $\mathcal{U} \gets \{ c \in \mathcal{C} : \delta(c; M) \ge m \}$
  \IF{$\mathcal{U} = \emptyset$}
    \STATE $\mathcal{U} \gets \mathcal{C}$ \COMMENT{no usable candidate: keep the closest one}
  \ENDIF
  \STATE $s_{t+1} \gets \arg\max_{c \in \mathcal{U}} \delta(c; M)$; \ $t \gets t + 1$
\ENDWHILE
\RETURN $s_1 s_2 \dots s_t$
\end{algorithmic}
\end{algorithm}

\subsection{Uncertainty-gated push}
\label{sec:push}

The push raises the per-draw success probability by tilting the model's word choices, while a candidate sentence is being
written, toward continuations that move the sentence deeper into the mask. Let $\ell \in \mathbb{R}^{|V|}$ be the
model's next-token logits at step $j$ of a candidate sentence whose first $j - 1$ tokens form the partial sentence
$c_{<j}$. Let $N$ be the nucleus (top-$p$ set, $p = 0.95$) of $\operatorname{softmax}(\ell / \tau)$ with $\tau = 0.8$,
and let $q$ be the model's sampling distribution: $\operatorname{softmax}(\ell / \tau)$ restricted to $N$ and
renormalised.

\paragraph{Entropy gate.} The strength of the push depends on the lexical uncertainty $H_j$ of
Eq.~\eqref{eq:token-entropy}, computed from the untempered distribution $\operatorname{softmax}(\ell)$:
\begin{equation}
  g(H_j) \;=\; \min\!\Big( 1, \, \max\!\Big( 0, \, \frac{H_j - H_{\text{lo}}}{H_{\text{hi}} - H_{\text{lo}}} \Big) \Big) ,
  \label{eq:gate}
\end{equation}
with $H_{\text{lo}} = 0.2$ and $H_{\text{hi}} = 2.0$ nats. At confident steps ($H_j \le H_{\text{lo}}$) the distribution
is left untouched (P2); entropy gating of this kind is also used by token-level watermarks for low-entropy text
\citep{lee2024sweet}.

\paragraph{Candidates.} Tokens outside $N$ keep probability zero, so the push never introduces a token the model would
not have sampled. We split $N$ into \emph{boundary} tokens
$N_{\mathcal{B}}$ (punctuation, whitespace, line breaks and end-of-text) and \emph{word} tokens $N_{\mathcal{W}}$. Only
the $L = 20$ most probable word tokens, $C \subseteq N_{\mathcal{W}}$, receive a bias.

\paragraph{Scores.} Each $v \in C$ is scored by the depth that the partial sentence would have with $v$ appended,
\begin{equation}
  u(v) \;=\; \min\!\big( \delta(c_{<j} \oplus v; M), \; \delta_{\max} \big) ,
  \label{eq:score}
\end{equation}
where $c_{<j} \oplus v$ is the decoded text of the partial sentence followed by $v$ and $\delta_{\max} = 0.05$. All $L$
strings are embedded in one batched encoder call. The cap makes the push switch itself off: once the partial sentence is
$\delta_{\max}$ deep inside the mask, all scores are equal and no token is favoured.

\paragraph{Bias.} The scores are standardised under the model's own probabilities. With
$\pi(v) = q(v) / \sum_{w \in C} q(w)$, $\mu = \sum_{v \in C} \pi(v)\, u(v)$ and
$\varsigma^2 = \sum_{v \in C} \pi(v)\, (u(v) - \mu)^2$, the bias of a token is
\begin{equation}
  \Delta(v) \;=\; \beta \, g(H_j) \, \operatorname{clip}\!\Big( \frac{u(v) - \mu}{\varsigma}, \, -1, \, 1 \Big)
  \quad (v \in C),
  \label{eq:bias}
\end{equation}
with $\beta = 1$; $\Delta(v) = 0$ for the other tokens and whenever $\varsigma = 0$. A token that would complete a 4-gram
already present in the generated text receives no positive bias, which prevents the push from rewarding repetition.

\paragraph{Mass-preserving tilt.} The bias is added to the logits of the word tokens only, and the word class is
renormalised to its original total probability:
\begin{equation}
  \tilde q(v) \;=\;
  \begin{cases}
    q(v), & v \in N_{\mathcal{B}}, \\[3pt]
    Q_{\mathcal{W}} \, \dfrac{q(v)\, e^{\Delta(v)/\tau}}{\sum_{w \in N_{\mathcal{W}}} q(w)\, e^{\Delta(w)/\tau}}, & v \in N_{\mathcal{W}},
  \end{cases}
  \qquad Q_{\mathcal{W}} = \sum_{w \in N_{\mathcal{W}}} q(w) ,
  \label{eq:tilt}
\end{equation}
and the next token is sampled from $\tilde q$. Writing $\tilde p_\theta(\cdot \mid x, s_{1:t}; M)$ for the resulting
distribution over sentences, the push has the following properties by construction:
\begin{itemize}
  \item \textbf{Gated by uncertainty.} It acts only at steps with $H_j > H_{\text{lo}}$, with full strength from
  $H_{\text{hi}}$ on, and stops once the sentence is $\delta_{\max}$ deep in the mask.
  \item \textbf{Bounded.} Before renormalisation, a word token's logit moves by at most $\beta$; the push only reweights
  tokens inside the model's nucleus, and sampling stays stochastic.
  \item \textbf{Boundary-preserving.} At every step the probability of a boundary token, and hence of ending the
  sentence or the answer, equals the model's own given the same prefix (P3).
\end{itemize}
The push is a reweighting of the model's sampling distribution, not a filter: the selection of
Section~\ref{sec:selection} still decides which candidate is kept. The two parts are complementary. The push raises the
chance that a single draw is usable at positions with lexical freedom, so the first block succeeds more often; the
bounded selection limits the cost at positions without semantic freedom, where neither resampling nor steering can move
the sentence into the mask.

\subsection{Sentence units}
\label{sec:units}

The mark is carried by the sentences the detector sees, so the generator must segment text exactly as the detector does.
A candidate sentence ends at sentence punctuation (optionally followed by closing quotes or brackets) followed by
whitespace, or at a line break. Three rules keep the generator's units identical to the detector's. First, line breaks
that a tokenizer emits as separate tokens are joined to the following sentence instead of forming a sentence of their
own, which would shift the chain of masks. Second, before each new sentence the model makes one ordinary draw of its next
token, and the answer ends if that draw is end-of-text; inside a sentence, end-of-text is blocked only at the first token.
Third, the first sentence is never empty: if the model ends the answer before writing one, it is sampled again.

\subsection{Detection}
\label{sec:detection}

Detection is SemStamp's, unchanged, and needs only the text and the key. The text is split into sentences
$s_1, \dots, s_n$, each sentence is embedded, and the detector counts the sentences $s_2, \dots, s_n$ whose region lies in
the mask of their predecessor, $\kappa = |\{ t \ge 2 : r(s_t) \in \mathcal{M}_k(r(s_{t-1})) \}|$. With $n' = n - 1$
scored sentences,
\begin{equation}
  z \;=\; \frac{\kappa - \gamma n'}{\sqrt{\gamma (1 - \gamma)\, n'}}
  \;=\; \sqrt{n'} \; \frac{\hat\pi - \gamma}{\sqrt{\gamma (1 - \gamma)}} ,
  \qquad \hat\pi = \kappa / n' ,
  \label{eq:z}
\end{equation}
and the text is flagged if $z$ exceeds a threshold calibrated on human-written text. For a fixed number of sentences,
the statistic grows linearly with the share $\hat\pi$ of sentences that lie in their mask. Raising $\hat\pi$ is
therefore the only lever a generator has on detection; SemStamp loses it at every sentence it gives up on, and \ours{}
is designed to keep it high while drawing at most $K$ candidates per sentence.
